\scanpage{8}
\mspage{6}
et pour $x,y,z\in\operatorname{Ob}X$, la composition
\[
\operatorname{Hom}_\Sigma(x,y)\times\unsure{H\text{--}|\operatorname{Hom}}_\Sigma(y,z)\overset{(u,v)\to vu}{\unsure{\text{---}|\longrightarrow}}{H\text{--}}_{\Sigma}(x,z),
\]
(où on met $\operatorname{Hom}_\Sigma$ au lieu de $\operatorname{Hom}_{X\Sigma^{-1}}$, pour alléger les notations).

\deleted{Je \unsure{note} d} Grâce à $\Sigma$, on définit une sous-catégorie pleine
\[
\underline{\mathrm{Ch}}_{\Sigma\deleted{\unsure{X}}}(X)\subset\underline{\mathrm{Ch}}(X),
\]
% slanted
\margin{L'ens de ses objets est noté $\mathrm{Ch}_\Sigma(X)$, ses \unsure{élémts|éléments} sont les $\Sigma$-chemins de $x$ à $y$}
par la condition sur $c\in\deleted{\unsure{\operatorname{Ob}}}\mathrm{Ch}(X)$
\[
\begin{aligned}
&\deleted{\text{Pour }\unsure{i\to i+1|i\ \text{pair}}}\\
&\text{Pour }i\text{ impair i.e.}\quad i-1\mathrel{\underset{\deleted{\unsure{\to}}}{\xleftarrow{u_i}}}i\quad (u_i\ \deleted{\unsure{\text{est dirigée}}}\ \inserted{\text{est dite}}\\
&\text{``rétrograde''}),\quad c(u_i)\in\Sigma,
\end{aligned}
\]
i.e. les \deleted{\unsure{tronçons}} \deleted{flèches} ``tronçons'' rétrogrades dans le chemin $c$ doivent être des $\in\Sigma$. On définit de même les sous-catégories pleines
\[
\underline{\mathrm{Ch}}_\Sigma(X;x,y)=\{c\in\underline{\mathrm{Ch}}(X;x,y)\mid c\in\mathrm{Ch}_\Sigma(X)\}\unsure{.}
\]
Ou encore la fibre en $(x,y)$ du foncteur source-but pour $\underline{\mathrm{Ch}}_\Sigma(X)$
\[
\underline{\mathrm{Ch}}_\Sigma(X)\longrightarrow X\times X.
\]

On définit \unsure{une} application canonique
\[
\mathrm{Ch}_\Sigma(X;x,y)\longrightarrow\operatorname{Hom}_\Sigma(x,y),
\]

% archive sheet number: 7
